% Generated draft: requires mathematical and source review.
\documentclass[11pt]{article}
\usepackage[margin=25mm]{geometry}
\usepackage{fontspec}
\setmainfont{Latin Modern Roman}
\usepackage{amsmath,amssymb,mathtools,unicode-math}
\setmathfont{Latin Modern Math}
\usepackage{enumitem,needspace,xurl,hyperref}
\setlength{\parindent}{0pt}
\setlength{\parskip}{5pt}
\setlength{\emergencystretch}{3em}
\newcommand{\N}{\mathbb{N}}
\newcommand{\Z}{\mathbb{Z}}
\newcommand{\Q}{\mathbb{Q}}
\newcommand{\R}{\mathbb{R}}
\newcommand{\C}{\mathbb{C}}
\newcommand{\F}{\mathbb{F}}
\newcommand{\A}{\mathbb{A}}
\newcommand{\PP}{\mathbb{P}}
\newcommand{\E}{\mathbb{E}}
\newcommand{\eps}{\varepsilon}
\newcommand{\dd}{\,\mathrm{d}}
\begin{document}
\section*{UnsolvedMath ID: 30003321}
\textbf{Source problem code:} catalogue entry
\subsection{Definitions and mathematical statement}
Let \(\mathbf{No}\) denote the ordered field of surreal numbers, with its usual addition, order, and simplicity relation. Write \(y\prec_s x\) when the sign expansion of \(y\) is a proper initial segment of the sign expansion of \(x\). A subclass \(S\subseteq\mathbf{No}\) is initial if
\[
 x\in S\ \text{and}\ y\prec_s x\quad\Longrightarrow\quad y\in S.
\]
An additive subgroup \(G\subseteq\mathbf{No}\), with the inherited order, is called discrete here if it has a least positive element:
\[
 \exists u\in G\;\bigl(u>0\ \land\ \forall g\in G\,(g>0\Rightarrow u\le g)\bigr).
\]
This convention excludes the trivial subgroup; including it separately does not affect the question, since it is itself initial and contained in \(\mathbf{Oz}\).

The omnific integers \(\mathbf{Oz}\) are the surreal numbers whose Conway normal form has only nonnegative exponents and an integer coefficient at exponent zero. Explicitly, if
\[
 x=\sum_{\alpha<\lambda}r_\alpha\omega^{a_\alpha},
\]
where \(\lambda\) is an ordinal, the exponents \(a_\alpha\in\mathbf{No}\) strictly decrease, and \(r_\alpha\in\R\setminus\{0\}\), then \(x\in\mathbf{Oz}\) precisely when every \(a_\alpha\ge0\) and any coefficient with \(a_\alpha=0\) is an integer. The empty sum represents zero.

The supplied problem asks: for every discrete initial additive subgroup \(G\subseteq\mathbf{No}\), does there exist an initial additive subgroup \(H\subseteq\mathbf{Oz}\) isomorphic to \(G\)? Here initiality of \(H\) is understood using the ambient surreal simplicity relation.

There are two precise readings of the unspecified word ``isomorphic.'' The abstract group reading requires a bijection \(\phi:G\to H\) satisfying
\[
 \phi(g_1+g_2)=\phi(g_1)+\phi(g_2)
 \qquad(g_1,g_2\in G).
\]
The ordered group reading additionally requires
\[
 g_1<g_2\quad\Longleftrightarrow\quad\phi(g_1)<\phi(g_2)
 \qquad(g_1,g_2\in G).
\]
The supplied record does not explicitly distinguish these readings. It also does not specify whether the quantified subgroups may be proper classes or must be sets. These scope and isomorphism qualifications require source review before a single definitive formalization is selected.
\subsection{Short English statement}
Does every discrete additive subgroup of the surreal numbers that is closed under simpler predecessors have an isomorphic initial copy inside the omnific integers? The source leaves preservation of order and the allowance of proper class subgroups unspecified.
\subsection{Sources}
\begin{itemize}
\item Local UnsolvedMath registry, supplied complete record with permanent ID \(30003321\), source problem code \texttt{catalogue entry}, title ``Discrete Initial Subgroups of the Surreal Numbers.'' No individual proposer is supplied; the proposed year is recorded as 2016.
\item Supplied source citation: \emph{Mini-Workshop: Surreal Numbers, Surreal Analysis, Hahn Fields and Derivations} (2017), Oberwolfach Reports, pp.\ 3313--3372, DOI \texttt{10.4171/owr/2016/60}. \url{https://doi.org/10.4171/owr/2016/60}. This is also the supplied literature source labelled ``original OWR report.''
\item Supplied literature source: \emph{Number systems with simplicity hierarchies II}. \url{https://arxiv.org/abs/1512.04001}.
\item Literature assessment dated exactly \texttt{2026-08-21}: the supplied curation labels the problem open and states that discrete initial subdomains have been classified, while the analogous question for arbitrary discrete initial subgroups remains open. This is supplied dated curation, not an independently verified literature assessment.
\item The registry labels the statement \texttt{corrected\_verified}, explaining that notation, typography, and sentence structure were repaired without changing the mathematical content. This supplied editorial label does not establish mathematical correctness or resolve the interpretive qualifications above.
\end{itemize}
\end{document}
